\documentclass[../mathNotesPreamble]{subfiles}

\providecommand{\relscalefact}{1.4}
\begin{document}
\relscale{\relscalefact}
  \section{3.3: Power Functions and Polynomial Functions}
  \begin{defn*}[Power Function]
    A \textbf{power function} is a function that can be represented in the form
      \[f(x)=kx^p\]
    where $k$ and $p$ are real numbers, and $k$ is known as the \textbf{coefficient}.
  \end{defn*}
  \begin{ex*}
    \TabPositions{35mm}
    Identify which of the following functions are power functions:
    \begin{extasks}[after-item-skip=\stretch{1}](2)
      \task $f(x)=1$ \tab Constant
      \task $f(x)=x$ \tab Identity
      \task $f(x)=x^2$ \tab Quadratic
      \task $f(x)=x^3$ \tab Cubic
      \task $f(x)=\dfrac{1}{x}$ \tab Reciprocal
      \task $f(x)=\dfrac{1}{x^2}$ \tab Reciprocal squared
      \task $f(x)=\sqrt{x}$ \tab Square root
      \task $f(x)=\sqrt[3]{x}$ \tab Cube root
      \task $f(x)=2^x$ \tab Exponential
    \end{extasks}
    \vspace*{\stretch{1}}
  \end{ex*}
  \pagebreak

  \noindent
  \textbf{Identifying End Behavior of Power Functions}

  \noindent
  Below are example graphs of power functions $f(x)=kx^n$:
  \begin{center}
    \begin{tikzpicture}[
      funcLabel/.style={circle, inner sep=0.5pt, fill=white,opacity=0.0, text opacity=1}]
      \begin{groupplot}[
        group style={group size=2 by 2,horizontal sep=3cm, vertical sep=4cm},
        axis on top,
        axis lines=center,
        axis line style={black,->},
        xmin=-2, xmax=2,
        enlargelimits={abs=0.5},
        xlabel=, ylabel=,
        clip=false,
        ]
        \nextgroupplot[ymin=-1, ymax=4, title={$n$ Even, $k>0$}]
          \addplot[<->] expression[domain=-2:2]{x^2}
            node[pos=0.95, right, funcLabel] {$x^2$};
          \addplot[<->, red!50!black] expression[domain=-1.414:1.414]{x^4}
            node[pos=0.95, right, funcLabel] {$x^4$};
          \addplot[<->, green!50!black] expression[domain=-1.26:1.26]{x^6}
            node[pos=0.95, left, funcLabel] {$x^6$};
        %
        \nextgroupplot[ymin=-4, ymax=4, title={$n$ Odd, $k>0$}]
          \addplot[<->] expression[domain=-1.587:1.587]{x^3}
            node[pos=0.95, right, funcLabel] {$x^3$};
          \addplot[<->, red!50!black] expression[domain=-1.320:1.320]{x^5}
            node[pos=1.01, above right, funcLabel] {$x^5$};
          \addplot[<->, green!50!black] expression[domain=-1.219:1.219]{x^7}
            node[pos=0.95, left, funcLabel] {$x^7$};
        %
        \nextgroupplot[ymin=-4, ymax=1, title={$n$ Even, $k<0$}]
          \addplot[<->] expression[domain=-2:2]{-x^2}
            node[pos=0.95, right, funcLabel] {$-x^2$};
          \addplot[<->, red!50!black] expression[domain=-1.414:1.414]{-x^4}
            node[pos=0.95, below, funcLabel] {$-x^4$};
          \addplot[<->, green!50!black] expression[domain=-1.26:1.26]{-x^6}
            node[pos=0.95, left, funcLabel] {$-x^6$};
        %
        \nextgroupplot[ymin=-4, ymax=4, title={$n$ Odd, $k<0$}]
          \addplot[<->] expression[domain=-1.587:1.587]{-x^3}
            node[pos=0.95, right, funcLabel] {$-x^3$};
          \addplot[<->, red!50!black] expression[domain=-1.320:1.320]{-x^5}
            node[pos=0.975, below, funcLabel] {$-x^5$};
          \addplot[<->, green!50!black] expression[domain=-1.219:1.219]{-x^7}
            node[pos=0.95, left, funcLabel] {$-x^7$};
      \end{groupplot}
    \end{tikzpicture}
  \end{center}
  \pagebreak

  \begin{ex*}
    Describe the end behavior of the graphs of the following power functions:
    \begin{extasks}[after-item-skip=\stretch{1}](2)
      \task $f(x)=x^8$
      \task $g(x)=-x^9$
    \end{extasks}
    \vspace*{\stretch{1}}
  \end{ex*}

  \begin{defn*}[Polynomial Functions]
    Let $n$ be a non-negative integer. A \textbf{polynomial function} is a function that can be written in the form
      \[f(x)=a_nx^n+a_{n-1}x^{n-1}+\dots+a_2x^2+a_1x+a_0\]
    Each $a_i$ is a coefficient and can be any real number, and $a_n\neq 0$.\\[\baselineskip]

    The \textbf{degree} of the polynomial is the highest power that the variable is raised to. The \textbf{leading coefficient} is the coefficient associated with the term containing the highest power of the variable.
  \end{defn*}
  \begin{ex*}
    Describe the end behavior of the graphs of the following polynomials:
    \hspace*{\stretch{1}}
    \textcolor{blue}{\underline{\href{https://www.desmos.com/calculator/lewqpko988}{Graphs}}}
    \begin{extasks}[after-item-skip=\stretch{1}](2)
      \task $f(x)=5x^{4}+2x^{3}-x-4$
      \task $g(x)=-2x^{6}-x^{5}+3x^{4}+x^{3}$
      \task $h(x)=3x^{5}-4x^{4}+2x^{2}+1$
      \task $j(x)=-6x^{3}+7x^{2}+3x+1$
    \end{extasks}
    \vspace*{\stretch{1}}
  \end{ex*}
  \pagebreak

  \begin{defn*}[Intercepts and Turning Points of Polynomial Functions]
    A \textbf{turning point} is a point at which the graph changes direction from increasing to decreasing, or decreasing to increasing.\\

     The $y$-intercept is the point at which the function has an input value of zero.\\

     The $x$-intercepts are the points at which the output value is zero.\\
     
     A polynomial of degree $n$ will have, at most, $n$ $x$-intercepts, and $n-1$ turning points.
  \end{defn*}
  
  \begin{thmBox*}[Complex Conjugate Root Theorem]
    If a polynomial $f(x)$ has real coefficients, and $a+bi$ is a root of $f(x)$, then $a-bi$ is also a root of $f(x)$.
  \end{thmBox*}
  \begin{ex*}
    Find the intercepts of 
      \[f(x)=x^4-4x^2-45\]
    \vspace*{\stretch{1}}
  \end{ex*}

  \pagebreak
\end{document}
